\section*{Chapitre \ref{ch:b3:seismology-imaging} --- Sismologie et imagerie de la Terre}

\begin{solution}{exo:b3:seismology-imaging:1}
(a) $\sin i_0 = pv_0/R = 500\times5.80/6371 = 0.4552$, donc $i_0 = \ang{27.1}$.
(b) À \qty{771}{km}, $r = \qty{5600}{km}$ et $r/v = 5600/11.066 =
\qty{506.1}{s.rad^{-1}}$~; à \qty{871}{km}, $r = \qty{5500}{km}$ et $r/v =
5500/11.245 = \qty{489.1}{s.rad^{-1}}$. Le rai atteint son \omterm{def:b3:seismology-imaging:ray-parameter}{point de retournement} lorsque $r/v = p = 500$~:
$z_t = 771 + 100\times(506.1 - 500)/(506.1 - 489.1) = 771 + 36 =
\qty{807}{km}$.
\end{solution}

\begin{solution}{exo:b3:seismology-imaging:2}
(a) $(455.5 - 369.6)/10 = \qty{8.59}{s}$ par degré, et
$8.59\times180/\pi = \qty{492.2}{s.rad^{-1}}$. (b) Un degré d'arc vaut
$6371\pi/180 = \qty{111.19}{km}$, donc le front d'onde balaie la surface à
$111.19/8.59 = \qty{12.94}{km.s^{-1}}$, bien plus rapide que toute roche à la
surface~: le rai arrive presque verticalement. (c) Au \omterm{def:b3:seismology-imaging:ray-parameter}{point de retournement} $v = r_t/p =
5450/492.2 = \qty{11.07}{km.s^{-1}}$, soit une vitesse d'onde P dans le manteau inférieur.
\end{solution}

\begin{solution}{exo:b3:seismology-imaging:3}
(a) $\mu_1 = \rho_1\beta_1^2 = 2700\times3500^2 = \qty{3.31e10}{Pa}$
(\qty{33.1}{GPa})~; $\mu_2 = 3300\times4500^2 = \qty{6.68e10}{Pa}$
(\qty{66.8}{GPa}). (b) $\lambda = cT = 3.79\times20 = \qty{75.8}{km}$, soit 2.17
fois l'épaisseur de la couche~: l'onde pénètre bien dans le demi-espace, ce qui
explique pourquoi sa vitesse est déjà supérieure à $\beta_1$. (c) $4000/3.384 = \qty{1182}{s}$ et
$4000/4.333 = \qty{923}{s}$~: les ondes de 20~s arrivent \qty{259}{s}, soit
\qty{4.3}{min}, après les ondes de 100~s.
\end{solution}

\begin{solution}{exo:b3:seismology-imaging:4}
(a) ${}_0S_2$~: $1/(53.89\times60) = \qty{3.093e-4}{Hz} = \qty{0.3093}{mHz}$~;
${}_0T_2$~: $1/(44.01\times60) = \qty{0.3787}{mHz}$. (b) $1440/53.89 = 26.7$
et $1440/44.01 = 32.7$ oscillations par jour. (c) La différence vaut
$\qty{0.0694}{mHz}$, dont l'inverse vaut \qty{14400}{s}~: il faut au moins
\qty{4.0}{h} d'enregistrement, en pratique bien davantage, car les pics sont élargis
par l'atténuation et dédoublés.
\end{solution}

\begin{solution}{exo:b3:seismology-imaging:5}
(a) La corde entre deux points de la sphère séparés par l'angle $\Delta$
au centre a pour longueur $2R\sin(\Delta/2)$, parcourue à la vitesse $v$.
(b) À \ang{30}~: $2\times6371\times\sin\ang{15}/10.0 = \qty{329.8}{s}$, contre
\qty{369.6}{s}~; à \ang{90}~: $2\times6371\times\sin\ang{45}/10.0 =
\qty{901.0}{s}$, contre \qty{779.7}{s}. (c) Le rai à \ang{30} atteint son \omterm{def:b3:seismology-imaging:ray-parameter}{point de retournement}
dans le manteau supérieur, où la vitesse réelle est d'environ 8 à \qty{10}{km.s^{-1}}, inférieure à la
valeur uniforme de \qty{10}{km.s^{-1}}~; le rai à \ang{90} passe l'essentiel de son trajet dans le
manteau inférieur, où la vitesse vaut de 11 à \qty{13.7}{km.s^{-1}}, et il s'y enfonce, si bien que la
Terre réelle est plus rapide.
\end{solution}

\begin{solution}{exo:b3:seismology-imaging:6}
(a) En élevant au carré, $(2 - \xi)^4 = 16(1 - \xi/3)(1 - \xi)$. Développons~:
$16 - 32\xi + 24\xi^2 - 8\xi^3 + \xi^4 = 16 - \tfrac{64}{3}\xi + \tfrac{16}{3}\xi^2$,
d'où $\xi^4 - 8\xi^3 + \tfrac{56}{3}\xi^2 - \tfrac{32}{3}\xi = 0$, soit
$\xi\,(3\xi^3 - 24\xi^2 + 56\xi - 32) = 0$, et
$3\xi^3 - 24\xi^2 + 56\xi - 32 = (\xi - 4)(3\xi^2 - 12\xi + 8)$. La racine
$\xi = 0$ n'est pas une onde. (b) Les racines du trinôme sont $2 \pm 2/\sqrt3$,
soit 3.155 et 0.8453~; seule la valeur $\xi = 0.8453$ appartient à $(0, 1)$ (et $\xi = 4$
n'y appartient pas), et elle vérifie l'équation non élevée au carré, les deux membres valant $4/3$.
Donc $c/\beta = \sqrt{0.8453} = 0.9194$. (c) $c = 0.9194\times3.46 =
\qty{3.181}{km.s^{-1}}$~; sur \qty{1000}{km}, $314.4 - 289.0 = \qty{25.3}{s}$
après l'onde S.
\end{solution}

\begin{solution}{exo:b3:seismology-imaging:7}
$\mathsf G$ a pour lignes $(1, x_i)$ et $\mathsf C_d^{-1} = \operatorname{diag}
(100, 100, 25, 25)$ (\unit{s^{-2}}). Les équations normales sont
\[
  \begin{pmatrix} 250 & 23750 \\ 23750 & \num{2812500} \end{pmatrix}
  \begin{pmatrix} a \\ b \end{pmatrix}
  = \begin{pmatrix} 3665 \\ 418125 \end{pmatrix},
\]
de déterminant \num{1.390625e8}. D'où $b = (250\times418125 -
23750\times3665)/\num{1.390625e8} = \qty{0.12575}{s.km^{-1}}$ et
$a = (2812500\times3665 - 23750\times418125)/\num{1.390625e8} =
\qty{2.713}{s}$. La covariance est $\mathsf N^{-1}$~: $\sigma_b =
\sqrt{250/\num{1.390625e8}} = \qty{0.00134}{s.km^{-1}}$, $\sigma_a =
\sqrt{2812500/\num{1.390625e8}} = \qty{0.142}{s}$. La vitesse est
$1/b = \qty{7.95}{km.s^{-1}}$, avec l'erreur $\sigma_b/b^2 =
\qty{0.08}{km.s^{-1}}$~: une vitesse propre au sommet du manteau. Les résidus valent $-0.001$,
$0.011$, $-0.076$ et $\qty{0.036}{s}$, et $\chi^2 = 0.19$ pour $4 - 2 = 2$
degrés de liberté~: une valeur plus faible que prévu, donc les erreurs annoncées sont
plutôt généreuses~; l'ajustement est acceptable.
\end{solution}

\begin{solution}{exo:b3:seismology-imaging:8}
(a) $\mathsf G = (30\;\;40)$ est une seule ligne~: $\mathsf G^{\mathsf T}\mathsf G$ est
$2\times2$ de rang un, donc singulière~; tout modèle vérifiant $30m_1 + 40m_2 = 0$ peut
être ajouté sans modifier la donnée. (b) $\mathsf G\mathsf G^{\mathsf T} =
900 + 1600 = 2500$, donc $\mathbf m = (30, 40)\times0.50/2500 =
(0.006, 0.008)\,\unit{s.km^{-1}}$. (c) $\mathsf R = \frac{1}{2500}
\begin{pmatrix} 900 & 1200 \\ 1200 & 1600 \end{pmatrix} =
\begin{pmatrix} 0.36 & 0.48 \\ 0.48 & 0.64 \end{pmatrix}$, de trace 1~: une
donnée ne résout qu'une combinaison des paramètres. (d) $\mathsf R\,(0.01, 0) =
(0.0036, 0.0048)\,\unit{s.km^{-1}}$~: l'anomalie se répartit le long du rai, davantage
dans la cellule où le rai est plus long.
\end{solution}

\begin{solution}{exo:b3:seismology-imaging:9}
(a) En linéarisant~: $\delta t = -(\delta v/v)L/v = -0.03\times200/9.0 =
\qty{-0.667}{s}$~; exactement, $L/(1.03v) - L/v = 200/9.27 - 200/9.0 =
\qty{-0.647}{s}$. (b) Le trajet dans le panneau plongeant est $100/\cos\ang{40} =
\qty{130.5}{km}$~: $-0.03\times130.5/9.0 = \qty{-0.435}{s}$ linéarisé et
$\qty{-0.422}{s}$ exactement. (c) La valeur exacte est la valeur linéarisée divisée par
$1.03$~: une erreur de \qty{3}{\percent}, du même ordre que $\delta v/v$ lui-même, comme
on l'attend d'une théorie au premier ordre.
\end{solution}

\begin{solution}{exo:b3:seismology-imaging:10}
(a) Le mode harmonique existe tant que sa partie dans le demi-espace décroît, $c < \beta_2$~; à sa
coupure $c = \beta_2$, donc $q_2 = 0$ et $\tan(kHq_1) = 0$, avec $kHq_1$ dans
$[n\pi, n\pi + \pi/2)$~: $kHq_1 = n\pi$. (b) Avec $k = \omega/\beta_2$ et
$q_1 = \sqrt{\beta_2^2/\beta_1^2 - 1} = \sqrt{\beta_2^2 - \beta_1^2}/\beta_1$,
$\omega_nH\sqrt{\beta_2^2 - \beta_1^2}/(\beta_1\beta_2) = n\pi$, ce qui est la
formule. (c) $\sqrt{4.5^2 - 3.5^2} = \sqrt8 = 2.828$~; $\omega_1 =
\pi\times3.5\times4.5/(35\times2.828) = \qty{0.4998}{rad.s^{-1}}$, soit une période de
$2\pi/0.4998 = \qty{12.57}{s}$~: à des périodes plus longues, seul le mode fondamental
existe. (d) Pour $n = 0$, la fréquence de coupure est nulle~: le mode fondamental
existe à toutes les périodes. Aux courtes périodes, $kH \to \infty$ tandis que $kHq_1$ reste
inférieur à $\pi/2$, donc $q_1 \to 0$ et $c \to \beta_1$.
\end{solution}

\begin{solution}{exo:b3:seismology-imaging:11}
(a) $\mathsf G^{\mathsf T}(\mathsf G\mathsf G^{\mathsf T} + \varepsilon^2\mathsf I)
= (\mathsf G^{\mathsf T}\mathsf G + \varepsilon^2\mathsf I)\mathsf G^{\mathsf T}$;
multiplier à gauche par $(\mathsf G^{\mathsf T}\mathsf G + \varepsilon^2
\mathsf I)^{-1}$ et à droite par $(\mathsf G\mathsf G^{\mathsf T} +
\varepsilon^2\mathsf I)^{-1}$, toutes deux inversibles pour $\varepsilon > 0$.
(b) Ici $\mathsf G\mathsf G^{\mathsf T} = \|\mathbf g\|^2 = 2500$, un scalaire,
donc $\mathsf R_\varepsilon = \mathsf G^{\mathsf T}\mathsf G/(2500 + \varepsilon^2)$,
de trace $2500/(2500 + \varepsilon^2)$. (c) La donnée prédite est
$\mathbf g\cdot\mathbf m_\varepsilon = 0.50\times2500/(2500 + \varepsilon^2)$
et le résidu $0.50\,\varepsilon^2/(2500 + \varepsilon^2)$. Pour
$\varepsilon^2 = 0$~: trace 1, résidu 0~; pour 2500~: trace 0.5, résidu
\qty{0.25}{s}~; pour \num{10000}~: trace 0.2, résidu \qty{0.40}{s}. Un amortissement plus fort
donne un modèle plus petit, une moins bonne résolution et un écart d'ajustement plus grand.
\end{solution}

\begin{solution}{exo:b3:seismology-imaging:12}
(a) D'après l'exercice~5, $T = (2R/v)\sin(\Delta/2)$, donc $\mathrm dT/\mathrm d\Delta
= (R/v)\cos(\Delta/2)$. Une corde formant l'angle $\Delta$ au centre est la
base d'un triangle isocèle d'angle au sommet $\Delta$, elle rencontre donc le rayon
à la source en $i_0 = (\pi - \Delta)/2$, et $p = R\sin i_0/v =
(R/v)\cos(\Delta/2)$~: la même. (b) Le \omterm{def:b3:seismology-imaging:ray-parameter}{point de retournement} est le milieu de la
corde, à la distance $R\cos(\Delta_1/2)$ du centre~; de façon équivalente
$r_1 = p_1v$. (c) Ici $p(\Delta)/p_1 = \cos(\Delta/2)/\cos(\Delta_1/2)$, et
la formule donne $r_1 = R\exp\bigl(-\ln(1/\cos(\Delta_1/2))\bigr) =
R\cos(\Delta_1/2)$, comme en (b)~; la vitesse vaut $r_1/p_1 = R\cos(\Delta_1/2)\big/
\bigl((R/v)\cos(\Delta_1/2)\bigr) = v$.
\end{solution}

\begin{solution}{pb:b3:seismology-imaging:1}
\textbf{1.} $i_0 = \arcsin 0.8 = \ang{53.13}$ et $p = 0.8/5.0 =
\qty{0.16}{s.km^{-1}}$.

\textbf{2.} Le rai tourne là où $v = 1/p = \qty{6.25}{km.s^{-1}}$, soit à
$z = (6.25 - 5.0)/0.05 = \qty{25}{km}$.

\textbf{3.} En dérivant $\sin i = pv$ le long du rai, $\cos i\,\mathrm di/
\mathrm ds = pg\,\mathrm dz/\mathrm ds = pg\cos i$, donc $\mathrm di/\mathrm ds =
pg$~: la direction tourne à vitesse angulaire constante, soit la courbure d'un cercle de rayon
$1/(pg) = 1/(0.16\times0.05) = \qty{125}{km}$. Au \omterm{def:b3:seismology-imaging:ray-parameter}{point de retournement}, le rai est
horizontal à \qty{25}{km} de profondeur, et le centre se trouve verticalement à \qty{125}{km}
au-dessus de lui~: à $z = \qty{-100}{km}$, soit $v_0/g$ au-dessus de la surface, là où la
vitesse extrapolée s'annulerait.

\textbf{4.} À la source, le rayon du cercle est perpendiculaire au rai,
il fait donc l'angle $i_0$ avec l'horizontale, et le centre se trouve
à une distance horizontale de $125\cos i_0$~; par symétrie
$X = 2\times125\times\cos i_0 = 2\times125\times0.6 = \qty{150}{km}$, ou
$X = 2\sqrt{1 - p^2v_0^2}/(pg)$.

\textbf{5.} $\mathrm dT = \mathrm ds/v$ avec $\mathrm ds = \mathrm di/(pg)$ et
$v = \sin i/p$~: $\mathrm dT = \mathrm di/(g\sin i)$. Alors
\[
  T = \frac2g\int_{i_0}^{\pi/2}\frac{\mathrm di}{\sin i} = \frac2g\ln\cot\frac{i_0}2 =
  \frac2g\operatorname{arccosh}\frac1{pv_0} .
\]
Ici $\cot(\ang{26.57}) = 2$, donc
$T = 40\ln2 = \qty{27.73}{s}$.

\textbf{6.} Pour $p = 0.15$~: $pv_0 = 0.75$, $\sqrt{1 - 0.5625} = 0.6614$,
$X = 2\times0.6614/0.0075 = \qty{176.4}{km}$ et $T = 40\operatorname{arccosh}
(1.3333) = 40\times0.7954 = \qty{31.81}{s}$. Puis $\Delta T/\Delta X =
4.08/26.4 = \qty{0.155}{s.km^{-1}}$, compris entre 0.15 et 0.16, comme
l'exige $\mathrm dT/\mathrm dX = p$.

\textbf{7.} $\mu_1 = 2700\times3500^2 = \qty{33.1}{GPa}$ et
$\mu_2 = 3300\times4500^2 = \qty{66.8}{GPa}$.

\textbf{8.} $k = 2\pi/(cT) = 2\pi/(3.99\times25) = \qty{0.06299}{rad.km^{-1}}$;
$q_1 = \sqrt{(3.99/3.5)^2 - 1} = 0.5474$; $q_2 = \sqrt{1 - (3.99/4.5)^2} = 0.4624$.

\textbf{9.} $\tan(kHq_1) = 66.825\times0.4624/(33.075\times0.5474) = 1.707$~; le
mode fondamental donne $kHq_1 = \arctan1.707 = 1.041$, donc
$H = 1.041/(0.06299\times0.5474) = \qty{30.2}{km}$.

\textbf{10.} À \qty{40}{s}~: $k = 2\pi/170 = 0.03696$, $q_1 = 0.6888$,
$q_2 = 0.3287$, $\tan(kHq_1) = 0.964$, $kHq_1 = 0.7671$ et
$H = 0.7671/(0.03696\times0.6888) = \qty{30.1}{km}$. Les deux résultats concordent
à la précision près des vitesses de phase~: la croûte a une épaisseur d'environ \qty{30}{km}.

\textbf{11.} $\omega = 0.2513$ et $\qty{0.1571}{rad.s^{-1}}$, $k = 0.06299$ et
$\qty{0.03696}{rad.km^{-1}}$: $U \approx 0.0942/0.02603 = \qty{3.62}{km.s^{-1}}$,
inférieure aux deux vitesses de phase, comme il se doit lorsque $c$ croît avec la
période.

\textbf{12.} $\lambda = 3.99\times25 = \qty{99.8}{km}$ et $4.25\times40 =
\qty{170}{km}$. Le mode fondamental ne sonde qu'une fraction de sa
longueur d'onde~: l'onde de 25-s reste surtout dans la croûte, celle de 40-s atteint nettement
dans le manteau, si bien que sa vitesse est plus proche de $\beta_2$. Deux périodes donnant le
même $H$ constitue une vérification de l'adéquation du modèle à deux couches.

\textbf{13.} Avec les colonnes pour les cellules 1 à 4 et les lignes A à F, en unités
de \qty{1000}{km} ($\sqrt2\times100/1000 = 0.1414$)~:
\[
  \mathsf G = \begin{pmatrix}
  0.1 & 0.1 & 0 & 0 \\ 0 & 0 & 0.1 & 0.1 \\ 0.1 & 0 & 0.1 & 0 \\
  0 & 0.1 & 0 & 0.1 \\ 0.1414 & 0 & 0 & 0.1414 \\ 0 & 0.1414 & 0.1414 & 0
  \end{pmatrix}.
\]

\textbf{14.} $\mathsf C_d^{-1} = \operatorname{diag}(400, 400, 400, 400, 100,
100)$. Un rai de ligne ou de colonne contribue $0.1^2\times400 = 4$, un rai diagonal
$0.1414^2\times100 = 2$. La cellule 1 est traversée par A, C (4 chacun) et E (2)~: 10. Les cellules
1 et 2 partagent le rai A~: 4~; les cellules 1 et 3 partagent C~: 4~; les cellules 1 et 4 partagent E~: 2~; et
ainsi de suite par symétrie, ce qui donne les quatre lignes énoncées.

\textbf{15.} Un rai de ligne ou de colonne contribue $0.1\times400\,\delta t = 40\,
\delta t$ à chacune de ses cellules, un rai diagonal $0.1414\times100\,\delta t =
14.14\,\delta t$. Cellule 1~: $40(0.16 - 0.36) + 14.14\times(-0.15) = -10.12$~; cellule
2~: $40(0.16 + 0.11) + 14.14\times(-0.27) = 6.98$~; cellule 3~: $40(-0.48 - 0.36) +
14.14\times(-0.27) = -37.42$~; cellule 4~: $40(-0.48 + 0.11) + 14.14\times(-0.15) =
-16.92$.

\textbf{16.} Pour $(1,1,1,1)$, chaque ligne de $\mathsf N$ a pour somme 20. Pour
$(1,-1,-1,1)$, la ligne 1 donne $10 - 4 - 4 + 2 = 4$, et les autres lignes suivent le même
motif. Pour $(1,1,-1,-1)$, $10 + 4 - 4 - 2 = 8$, et pour $(1,-1,1,-1)$,
$10 - 4 + 4 - 2 = 8$. Les vecteurs sont orthogonaux de norme 2, donc
$\mathsf N^{-1} = \sum_k \mathbf u_k\mathbf u_k^{\mathsf T}/(4\lambda_k)$~; son
élément $(1,1)$ vaut $\tfrac14(\tfrac1{20} + \tfrac14 + \tfrac18 + \tfrac18) =
\tfrac{11}{80}$, son élément $(1,2)$ vaut $\tfrac14(\tfrac1{20} - \tfrac14 + \tfrac18
- \tfrac18) = -\tfrac4{80}$, son élément $(1,4)$ vaut $\tfrac14(\tfrac1{20} +
\tfrac14 - \tfrac18 - \tfrac18) = \tfrac1{80}$, et ainsi de suite.

\textbf{17.} $\hat{\mathbf m} = \mathsf N^{-1}\mathbf b$~:
$\hat m_1 = (11\times(-10.121) - 4\times6.982 - 4\times(-37.418) - 16.921)/80
= -0.081$~; de même $\hat m_2 = 1.844$, $\hat m_3 = -3.706$ et
$\hat m_4 = \qty{-0.931}{ms.km^{-1}}$.

\textbf{18.} $\delta v/v_0 = -v_0\,\delta s = -8.0\times10^{-3}\,\hat m$, soit en
pourcentage $-0.8\,\hat m$~: cellule 1 $+0.07$, cellule 2 $-1.48$, cellule 3 $+2.96$, cellule 4
$\qty{+0.75}{\percent}$.

\textbf{19.} Chaque élément diagonal de $\mathsf N^{-1}$ vaut $11/80$~: chaque cellule a
l'erreur type $\sqrt{0.1375} = \qty{0.371}{ms.km^{-1}}$, soit
$\qty{0.30}{\percent}$. La corrélation des cellules 1 et 2 vaut $-4/11 = -0.36$~:
elles partagent le rai A, et une lenteur attribuée à l'une peut être compensée par l'autre.
Les cellules 1 et 4, qui ne partagent que le rai diagonal E, moins précis, sont faiblement et
positivement corrélées, $+1/11 = +0.09$.

\textbf{20.} Les résidus prédits valent $0.176$, $-0.464$, $-0.379$, $0.091$,
$-0.143$ et $\qty{-0.263}{s}$~; les résidus après ajustement valent $-0.016$,
$-0.016$, $0.019$, $0.019$, $-0.007$ et $\qty{-0.007}{s}$, et
$\chi^2 = 0.50$ pour $6 - 4 = 2$ degrés de liberté~: compatible avec les erreurs
annoncées, voire un peu faible.

\textbf{21.} Les vecteurs propres sont ceux de $\mathsf N$ et les valeurs propres
$\lambda/(\lambda + 4)$~: $20/24 = 5/6$, $4/8 = 1/2$, $8/12 = 2/3$ deux fois. Comme
à la question 16, $R_{11} = \tfrac14(\tfrac56 + \tfrac12 + \tfrac23 + \tfrac23) =
\tfrac23$~; $R_{12} = \tfrac14(\tfrac56 - \tfrac12 + \tfrac23 - \tfrac23) =
\tfrac1{12}$~; $R_{13} = \tfrac14(\tfrac56 - \tfrac12 - \tfrac23 + \tfrac23) =
\tfrac1{12}$~; $R_{14} = \tfrac14(\tfrac56 + \tfrac12 - \tfrac23 - \tfrac23) = 0$.

\textbf{22.} On projette $\mathbf b$ sur les vecteurs propres~: $\mathbf b\cdot\mathbf u_k/4
= -14.37$, $0.849$, $12.80$ et $-9.40$ pour les valeurs propres 20, 4, 8 et 8.
L'estimation amortie les divise par $\lambda + 4$ au lieu de $\lambda$~; pour
la cellule 3 (composantes $1, -1, -1, 1$),
$m_3 = -14.37/24 - 0.849/8 - 12.80/12 - 9.40/12 = \qty{-2.555}{ms.km^{-1}}$,
soit une anomalie de vitesse de $\qty{+2.04}{\percent}$. L'amortissement ramène chaque composante
vers zéro, surtout les moins bien contraintes.

\textbf{23.} Le modèle non amorti, qui ajuste les données à leurs erreurs près
($\chi^2 = 0.50$ pour 2 degrés de liberté), alors que le modèle amorti présente
un écart d'ajustement important ($15.8$). L'amortissement n'est nécessaire que lorsque $\mathsf N$ est singulière ou
presque~; ici sa plus petite valeur propre, 4, n'est pas petite, l'estimation est stable et
ses erreurs sont modérées.

\textbf{24.} La ligne 3 est $(1/12, 0, 2/3, 1/12)$~: l'estimation amortie de la cellule 3 vaut
les deux tiers de sa vraie valeur plus un douzième de celles des cellules 1 et 4, ses voisines
le long du rai C et du rai B, et rien de la cellule 2.

\textbf{25.} La cellule 3, la cellule du bas à gauche, est rapide
d'environ \qty{3}{\percent}~; la cellule 2 est lente. Un panneau plongeant froid et dense de lithosphère océanique
qui plonge sous le continent est plus rapide que le manteau environnant, donc la
cellule rapide, profonde et du côté océanique de K, est le panneau plongeant~; la cellule lente au-dessus
d'elle, de l'autre côté, se trouve là où le coin mantellique est chaud.

\textbf{26.} D'après les questions 18 et 19, l'anomalie de vitesse P de la cellule du panneau plongeant est
\textbf{\qty[text-series-to-math=true, uncertainty-mode=separate]{3.0 \pm 0.3}{\percent}}.
\end{solution}
